\subsection{Application to the Brauer Group}

In this subsection, K is a field, and L is a Galois algebra over K with action \(\lambda: G \to \operatorname{Aut}_{K}(L)\) . We denote the degree of L over K by n.~We have \(n = \operatorname{Card}(G)\) .

THEOREM 3. --- Let \(\mathcal{E} = (\Gamma, \iota, \pi)\) be a \(\tau\) -extension of G by L*. The algebra \(\mathbf{A}[\mathcal{E}; \mathrm{L}]\) is central simple and of degree \(n^2\) over K. Moreover, the canonical homomorphism \(u\) from L to \(\mathbf{A}[\mathcal{E}; \mathrm{L}]\) is an isomorphism from L to a maximal commutative subalgebra of \(\mathbf{A}[\mathcal{E}; \mathrm{L}]\) .

Lemma 12. --- a) There does not exist any ideal of L, other than \(\{0\}\) and L that is invariant under G.

\begin{enumerate}
\def\labelenumi{\alph{enumi})}
\setcounter{enumi}{1}
\tightlist
\item
  Let g be an element of G distinct from the identity element, and let \(\mathfrak{a}_{g}\) be the ideal of L generated by the elements of the form \(x - \lambda(g)x\) , where x runs through L. We have \(\mathfrak{a}_{g} = L\) .
\end{enumerate}

Let \(\mathscr{S}\) be the set of maximal ideals of L; for any subset S of \(\mathscr{S}\), set \(\mathfrak{a}(S) = \bigcap_{\mathfrak{m} \in S} \mathfrak{m}\) . Since the ring L is commutative and semisimple (VIII, p.~310, Remark 6), the map \(W \mapsto \mathfrak{a}(W)\) is a bijection from the set of subsets of \(\mathscr{S}\) to the set of ideals of L (VIII, p.~142, Proposition 9). The ideal \(\mathfrak{a}(S)\) is invariant under G if and only if S is. Since G acts transitively on \(\mathscr{S}\) (VIII, p.~308, Theorem 2, vi)), the only subsets of \(\mathscr{S}\) invariant under G are \(\varnothing\) and \(\mathscr{S}\). Since \(\mathfrak{a}(\varnothing) = L\) and \(\mathfrak{a}(\mathscr{S}) = \{0\}\) , we have proved assertion a).

Let us prove b) by contradiction; suppose that we have \(\mathfrak{a}_g\neq \mathrm{L}\) . Let \(\mathfrak{m}\) be a maximal ideal of L containing \(\mathfrak{a}_g\) . We then have \(\lambda (g)x\equiv x\bmod \mathfrak{m}\) for every x in L. In particular, m is invariant under g, and \(\lambda(g)\) is induced by the identity on the field L/m. By property (vi) in Theorem 2 of VIII, p.~308, the element g of \(G_{m}\) is trivial, which contradicts the assumption of b).

Let us now prove Theorem 3. The vector space \(\mathbf{A}[\mathcal{E};\mathrm{L}]\) has finite dimension \(\operatorname{Card}(\mathrm{G})[\mathrm{L}:\mathrm{K}] = n^2\) over \(\mathrm{K}\) , by VIII, p.~316 and the assumed equality \(\operatorname{Card}(\mathrm{G}) = [\mathrm{L}:\mathrm{K}] = n\) .

Let \(\mathfrak{a}\) be a nonzero two-sided ideal of the algebra \(\mathbf{A}[\mathcal{E};\mathrm{L}]\) . We use the notation of VIII, p.~316. Every element \(a\) of \(\mathbf{A}[\mathcal{E};\mathrm{L}]\) can be written uniquely as \(a = \sum_{g\in \mathrm{G}}a_g\varepsilon_g\) with \(a_g\in \mathrm{L}\) for every \(g\in \mathrm{G}\) ; we denote by \(\Phi (a)\) the set of elements \(g\) of \(\mathrm{G}\) such that \(a_g\neq 0\) . By formula (32) of VIII, p.~316, we have the relation

\[
\varepsilon_ {g} \varepsilon_ {g ^ {\prime}} = c _ {\sigma} (g, g ^ {\prime}) \varepsilon_ {g g ^ {\prime}}\tag{38}
\]

for all \(g, g' \in \mathrm{G}\) and consequently

\[
\Phi (a \varepsilon_ {g}) = \Phi (a). g\tag{39}
\]

for every \(g\in \mathbf{G}\) and \(a\in \mathbf{A}[\mathcal{E};\mathrm{L}]\) .

Let \(a\) be a nonzero element of \(\mathfrak{a}\) for which \(\Phi(a)\) is minimal for the inclusion; by formula (39), if necessary replacing \(a\) with an element of the form \(a\varepsilon_{g^{-1}}\) with \(g\in\Phi(a)\) , we may assume that we have \(e\in\Phi(a)\) . Let \(s\) be an element of \(\Phi(a)\) distinct from \(e\) . By Lemma 12, b), there exists an element \(x\) of L such that \(a_s(x-\lambda(s)\cdot x)\neq0\) . But we have the relation

\[
x a - a x = \sum_ {g} a _ {g} (x - \lambda (g) \cdot x) \varepsilon_ {g},\tag{40}
\]

where the sum is taken over the elements g of \(\Phi(a)\) distinct from e. Because of the minimality of \(\Phi(a)\) , we then have xa - ax = 0, but this contradicts the assumption on x. We have, therefore, proved by contradiction that \(\Phi(a)\) contains only the identity element e of G, so we have \(a \in L\) .

Hence \(L \cap a\) is an ideal of L not reduced to \(\{0\}\) . Moreover, for every x in L, we have \(\varepsilon_{g}x\varepsilon_{g}^{-1} = \lambda(g) \cdot x\) , so \(L \cap a\) is invariant under G. By Lemma 12, we therefore have \(L \cap a = L\) , that is, \(L \subset a\) . Since L contains the identity element of \(A[E;L]\) , we have \(a = A[E;L]\) .

Since the algebra \(A[E;L]\) has nonzero finite dimension over K and its only two-sided ideals are \(\{0\}\) and \(A[E;L]\) , it is simple.

Let us prove that every element \(a\) of \(\mathbf{A}[\mathcal{E};\mathrm{L}]\) that commutes with L belongs to L. We write \(a\) as \(\sum_{g\in \mathrm{G}}a_g\varepsilon_g\) with coefficients \(a_{g}\) in L. For every \(x\) in L, we have \(xa - ax = 0\) , and relation (40) shows that we have

\[
a _ {g} (x - \lambda (g) \cdot x) = 0\tag{41}
\]

for every \(g \in \mathrm{G}\) and \(x \in \mathrm{L}\) . By Lemma 12, we therefore have \(a_g = 0\) for \(g \neq e\) ; consequently, \(a \in \mathrm{L}\) .

Let us, finally, determine the center of \(\mathbf{A}[\mathcal{E};\mathrm{L}]\) . If \(z\) belongs to the center, then it commutes with \(\mathrm{L}\) , so belongs to \(\mathrm{L}\) . But we then have

\[
0 = \varepsilon_ {g} z - z \varepsilon_ {g} = (\lambda (g) \cdot z - z) \varepsilon_ {g}
\]

for every g of G. Hence z is invariant under the automorphism group \(\lambda(\mathrm{G})\) of L. We therefore have \(z \in K\) by assumption.

THEOREM 4. --- Let A be a central simple algebra of finite degree over K, and let L be a maximal commutative subalgebra of A. Then there exists a \(\tau\) -extension \(\mathcal{E}\) of G by \(\mathrm{L}^*\) such that A is isomorphic to \(\mathbf{A}[\mathcal{E};\mathrm{L}]\) .

Without loss of generality, we may assume that L is a maximal commutative subalgebra of A. Let \(\Gamma\) be the multiplicative group consisting of the invertible elements \(\gamma\) of A such that there exists a g in G with

\[
\gamma x \gamma^ {- 1} = \lambda (g). x\tag{42}
\]

for every element x of L. If \(\gamma \in \Gamma\) is given, then the element g satisfying this relation is unique; we denote it by \(\pi(\gamma)\) . It is immediate that \(\pi\) is a homomorphism from \(\Gamma\) to G and that its kernel is equal to \(L^{*}\) . The surjectivity of \(\pi\) follows from the Skolem--Noether theorem (VIII, p.~263, Corollary).

If we denote the canonical injection of \(L^{*}\) into \(\Gamma\) by \(\iota\) , then it follows from the constructions that \(\mathcal{E} = (\Gamma, \iota, \pi)\) is a \(\tau\) -extension of G by \(L^{*}\) . Let

\[
u: \mathrm{L} \to \mathbf {A} [ \mathcal {E}; \mathrm{L} ] \quad \text { and } \quad v: \Gamma \to \mathbf {A} [ \mathcal {E}; \mathrm{L} ]
\]

be the canonical homomorphisms. By the universal property of \(\mathbf{A}[\mathcal{E};\mathrm{L}]\) (VIII, p.~315, Proposition 12), there exists a unique algebra homomorphism \(f:\mathbf{A}[\mathcal{E};\mathrm{L}]\to \mathrm{A}\) such that \(f\circ u(x) = x\) and \(f\circ v(\gamma) = \gamma\) for \(x\in \mathrm{L}\) and \(\gamma \in \Gamma\) . Since the algebra \(\mathbf{A}[\mathcal{E};\mathrm{L}]\) is simple, the homomorphism \(f\) is injective. Now, the algebra \(\mathbf{A}[\mathcal{E};\mathrm{L}]\) has degree \(n^2\) over \(\mathbf{K}\) and so does \(\mathrm{A}\) because \(\mathrm{L}\) is a maximal commutative semisimple subalgebra of \(\mathrm{A}\) and we have \(n = [\mathrm{L}:\mathrm{K}]\) (VIII, p.~262, Proposition 3, (ii)). Hence \(f\) is bijective.

DEFINITION 2. --- Let \(\mathscr{S}\) be the set of maximal ideals of L. We define

\[
\operatorname{Br} (\mathrm{L} / \mathrm{K}) = \bigcap_ {\mathfrak {m} \in \mathscr {S}} \operatorname{Ker} (r _ {(\mathrm{L} / \mathfrak {m}) / \mathrm{K}}),
\]

where \(r_{(L/\mathfrak{m})/K} : \mathrm{Br}(K) \to \mathrm{Br}(L/\mathfrak{m})\) is the extension of scalars homomorphism (VIII, p.~281).

THEOREM 5. --- There exists a group isomorphism

\[
\Psi : \mathrm{Ex} _ {\tau} (\mathrm{G}, \mathrm{L} ^ {*}) \longrightarrow \mathrm{Br} (\mathrm{L} / \mathrm{K})
\]

that sends the class of a \(\tau\) -extension \(\mathcal{E}\) of G by \(\mathrm{L}^*\) to the class in \(\mathrm{Br}(\mathrm{L} / \mathrm{K})\) of the algebra \(\mathbf{A}[\mathcal{E};\mathrm{L}]\) .

To define \(\Psi\) and verify that it is a bijection, we must establish the following points:

\begin{enumerate}
\def\labelenumi{\alph{enumi})}
\item
  If \(\mathcal{E}\) and \(\mathcal{E}'\) are isomorphic \(\tau\) -extensions of \(G\) by \(L^*\) , then the algebras \(\mathbf{A}[\mathcal{E}; L]\) and \(\mathbf{A}[\mathcal{E}'; L]\) are isomorphic.
\item
  Conversely, if the algebras \(\mathbf{A}[\mathcal{E};\mathrm{L}]\) and \(\mathbf{A}[\mathcal{E}';\mathrm{L}]\) are isomorphic, then the \(\tau\) -extensions \(\mathcal{E}\) and \(\mathcal{E}'\) of \(\mathrm{G}\) by \(\mathrm{L}^*\) are isomorphic.
\item
  In every class in \(\mathrm{Br(L / K)}\) , there is an algebra \(E\) containing \(L\) as a maximal commutative subalgebra.
\item
  If \(\mathrm{E}\) is a central simple algebra of finite degree over \(\mathrm{K}\) containing \(\mathrm{L}\) as a maximal commutative subalgebra, then there exists a \(\tau\) -extension \(\mathcal{E}\) of \(\mathrm{G}\) by \(\mathrm{L}^*\) such that \(\mathrm{E}\) is isomorphic to \(\mathbf{A}[\mathcal{E};\mathrm{L}]\) .
\end{enumerate}

Assertion a) follows from the construction of the cross product; b) follows from VIII, p.~263, Corollary. Assertion c) follows from Proposition 5 of VIII, p.~281, and d) is simply Theorem 4 above.

It remains to verify that \(\Psi\) is a group homomorphism; for this, it suffices to prove that if \(\mathcal{E}_1 = (\Gamma_1,\iota_1,\pi_1)\) and \(\mathcal{E}_2 = (\Gamma_2,\iota_2,\pi_2)\) are \(\tau\) -extensions, then the algebra \(\mathbf{A}[\mathcal{E}_1\mathcal{E}_2;\mathrm{L}]\) is equivalent to the algebra \(\mathbf{A}[\mathcal{E}_1;\mathrm{L}]\otimes \mathbf{A}[\mathcal{E}_2;\mathrm{L}]\) . We denote the product extension \(\mathcal{E}_1\mathcal{E}_2\) by \(\mathcal{E} = (\Gamma ,\iota ,\pi)\) . The group \(\Gamma\) is isomorphic to the cokernel of the homomorphism \(\rho\) from \(\mathrm{L}^*\) to the fiber product \(\Gamma_1\times_{\mathrm{G}}\Gamma_2\) that sends \(\mu\) to \((\iota_1(\mu),\iota_2(\mu)^{-1})\) . For \(i\in \{1,2\}\) , write \(\mathrm{A}_i = \mathbf{A}[\mathcal{E}_i;\mathrm{L}]\) . Let \(u_{i}:\mathrm{L}\to \mathrm{A}_{i}\) and \(v_{i}:\Gamma_{i}\rightarrow \mathrm{A}_{i}^{*}\) be the canonical homomorphisms. We identify L with its images by the homomorphisms \(u_{i}\) , which make L into maximal commutative subalgebras of the \(\mathrm{A}_i\) . We denote by \(\mathrm{V}_i\) the L-vector space defined by left multiplication in \(\mathrm{A}_i\) . The ring \(\mathrm{L}\otimes_{\mathrm{K}}\mathrm{A}_i^{\circ}\) acts on \(\mathrm{V}_i\) . Since it is simple and has dimension \(n^2\) over L, we obtain an isomorphism from \(\mathrm{L}\otimes_{\mathrm{K}}\mathrm{A}_i^{\circ}\) to \(\operatorname {End}_{\mathrm{L}}(\mathrm{V}_{i})\) . Hence the ring \(\mathrm{L}\otimes_{\mathrm{K}}\mathrm{A}_1^{\circ}\otimes_{\mathrm{K}}\mathrm{A}_2^{\circ}\) , which we can identify with \((\mathrm{L}\otimes_{\mathrm{K}}\mathrm{A}_1^{\circ})\otimes_{\mathrm{L}}(\mathrm{L}\otimes_{\mathrm{K}}\mathrm{A}_2^{\circ})\) , is isomorphic to \(\operatorname {End}_{\mathrm{L}}(\mathrm{V}_{1}\otimes_{\mathrm{L}}\mathrm{V}_{2})\) . Set \(\mathrm{C} = \operatorname {End}_{\mathrm{A}_1^{\circ}\otimes_{\mathrm{K}}\mathrm{A}_2^{\circ}}(\mathrm{V}_1\otimes_\mathrm{L}\mathrm{V}_2)\) . By Lemma 3 of VIII, p.~282, the ring C is similar to \(\mathrm{A}_1\otimes_{\mathrm{K}}\mathrm{A}_2\) , and \(\mathrm{L}\otimes 1\otimes 1\) is a maximal commutative subalgebra of C. For every pair \((\gamma_{1},\gamma_{2})\in \Gamma_{1}\times \Gamma_{2}\) satisfying \(\pi_1(\gamma_1) = \pi_2(\gamma_2)\) , we denote by \(w(\gamma_1,\gamma_2)\)

the unique \(\lambda (\pi_1(\gamma_1))\) -semilinear (II, §1, No.~13, p.~223) endomorphism such that

\[
w (\gamma_ {1}, \gamma_ {2}) (x _ {1} \otimes x _ {2}) = v _ {1} (\gamma_ {1}) x _ {1} \otimes v _ {2} (\gamma_ {2}) x _ {2}
\]

for \(x_{1} \in V_{1}\) and \(x_{2} \in V_{2}\) . We have \(w(\gamma_{1}, \gamma_{2}) \in \mathrm{C}^{*}\) , and w is a group homomorphism from the fiber product \(\Gamma_{1} \times_{G} \Gamma_{2}\) to \(C^{*}\) . This homomorphism is trivial on the image of \(\rho\) and induces a homomorphism v from \(\Gamma\) to \(C^{*}\) . Denote by \(u : L \to C\) the morphism given by \(u : \ell \mapsto \ell \otimes 1 \otimes 1\) . We can verify the relations

\[
u (\alpha) = v (\iota (\alpha)) \quad \text { and } \quad u (^ {\gamma} \beta) = v (\gamma) u (\beta) v (\gamma) ^ {- 1}
\]

for \(\alpha \in \mathrm{L}^*\) , \(\beta \in \mathrm{L}\) , and \(\gamma \in \Gamma\) . Proposition 12 of VIII, p.~315 gives a homomorphism \(f\) from the algebra \(\mathbf{A}[\mathcal{E};\mathrm{L}]\) to \(\mathrm{C}\) . Since the algebra \(\mathbf{A}[\mathcal{E};\mathrm{L}]\) is simple, the homomorphism \(f\) is injective. The algebras \(\mathrm{C}\) and \(\mathbf{A}[\mathcal{E};\mathrm{L}]\) have the same dimension over \(\mathrm{K}\) , so \(f\) is an isomorphism.

Remarks. --- 1) If L is an étale algebra over K and G is the automorphism group of L, then it is not always true that the algebra \(A[E;L]\) is central simple (for example, we can take \(L = K^{n}\) and \(G = S_{n}\) ).

\begin{enumerate}
\def\labelenumi{\arabic{enumi})}
\setcounter{enumi}{1}
\tightlist
\item
  We can calculate a 2-cocycle c associated with an algebra A split by a finite Galois extension L with group G as follows. First, there exists a K-homomorphism \(\varphi: A \to \mathbf{M}_{m}(L)\) , where \([A:K] = m^{2}\) . For \(g \in G\) , let \(\varphi^{g}\) be the homomorphism from A to \(\mathbf{M}_{m}(L)\) given by \(a \mapsto \varphi(g^{-1}ag)\) . By the Skolem--Noether theorem (VIII, p.~256, Theorem 3), for every \(g \in G\) , there exists an element \(u_{g}\) of \(\mathbf{GL}_{m}(L)\) such that
\end{enumerate}

\[
\varphi^ {g} (a) = u _ {g} \varphi (a) u _ {g} ^ {- 1}
\]

for \(a\in \mathrm{A}\) . We then set

\[
c (g, g ^ {\prime}) = u _ {g} u _ {g ^ {\prime}} u _ {g g ^ {\prime}} ^ {- 1}.
\]

We can also define an extension of G by \(L^{*}\) using \(\varphi\) : we consider the group \(\Gamma \subset \mathbf{GL}_{m}(L)\) consisting of the \(\gamma\) for which there exists a \(g \in G\) with

\[
\varphi^ {g} (a) = \gamma \varphi (a) \gamma^ {- 1}
\]

for every \(a \in \mathrm{A}\) . The class of this extension is the inverse image by \(\Psi\) of the class of \(\mathrm{A}\) in \(\mathrm{Br}(\mathrm{L} / \mathrm{K})\) .

COROLLARY. --- The mapping \(\Phi_{L/K} = \Theta \circ \Psi^{-1}\) defines a group isomorphism from \(\mathrm{Br}(\mathrm{L}/\mathrm{K})\) to \(\mathrm{H}^{2}(\mathrm{G}, \mathrm{L}^{*})\) .

Let \(K'\) be an extension of K and \(\varphi: K' \to L\) a morphism of K-algebras. The set H of elements h of G such that \(\lambda(h) \circ \varphi = \varphi\) is a subgroup of G, and the \(K'\) -algebra L endowed with the restriction of \(\lambda\) to H is a Galois algebra over \(K'\) .

PROPOSITION 14. --- The following diagram commutes:

\[
\begin{array}{c} \operatorname{Br} (\mathrm{L} / \mathrm{K}) \xrightarrow {\Phi_ {\mathrm{L} / \mathrm{K}}} \operatorname{H} ^ {2} (\operatorname{G}, \operatorname{L} ^ {*}) \\ \Biggl \downarrow^ {r _ {\mathrm{K} ^ {\prime} / \mathrm{K}}} \qquad \qquad \qquad \qquad \Biggl \downarrow^ {\operatorname{Res} _ {\mathrm{H}} ^ {\mathrm{G}}} \\ \operatorname{Br} (\mathrm{L} / \mathrm{K} ^ {\prime}) \xrightarrow {\Phi_ {\mathrm{L} / \mathrm{K} ^ {\prime}}} \operatorname{H} ^ {2} (\operatorname{H}, \operatorname{L} ^ {*})  . \end{array}
\]

This follows from Propositions 7 of VIII, p.~299 and 13 of VIII, p.~317.

